\begin{proposition}[Color-diagonal census and inverse-reversal]
Let $\Omega'$ be the $21{,}456$-element colored Nielsen component for
$(2A,2A,2A,2A,3A)$.  Let $q_1,q_2,q_3$ be the three adjacent
same-color Hurwitz generators and let $m=q_4^2$ be the adjacent
mixed-color collision operator.  Then
\[
\operatorname{ct}(q_i)=1^{24}2^{292}3^{796}4^{1248}5^{980}6^{1428}
\qquad (i=1,2,3),
\]
and
\[
\operatorname{ct}(m)=1^{18}2^{468}3^{582}4^{984}5^{900}
7^{840}8^{192}11^{264}.
\]
Hence each same-color width-$5$ locus contains exactly
$5\cdot980=4900$ states.  Among the $18$ fixed states of $m$, the
$q_1$- and $q_2$-width distributions are
\[
2^2,\qquad 3^6,\qquad 5^{10},
\]
so each of the two corresponding color-diagonal cells has exactly ten
states.  The $q_3$-width distribution on $\operatorname{Fix}(m)$ is
$4^5 6^{13}$, so that cell is empty.

Define the color-preserving inverse-reversal involution by
\[
c(g_1,g_2,g_3,g_4,g_5)
=(g_4^{-1},g_3^{-1},g_2^{-1},g_1^{-1},g_5^{-1}).
\]
It acts on all of $\Omega'$, satisfies $c^2=1$, and has $42$ fixed
states.  It conjugates
\[
cq_1c=q_3^{-1},\qquad cq_2c=q_2^{-1},\qquad cq_3c=q_1^{-1}.
\]
A chosen ten-state diagonal cell is therefore carried to a distinct
conjugate ten-state cell.  After including all six same-color pair
collisions and all four mixed-color pair collisions, the complete
color-diagonal set has $60$ states.  It is $c$-stable and $c$ acts on
it without fixed points, as $30$ transpositions.
\end{proposition}

\begin{proof}
Canonical simultaneous relabeling reconstructs the complete colored
orbit from the printed five-tuple.  Direct cycle decomposition gives
the two displayed passports.  Intersecting the fixed set of each mixed
collision operator with the width-$5$ states of each same-color
collision operator gives twelve empty cells, four cells of size four,
and eight cells of size ten; their union contains $60$ states.  Direct
application of $c$ to the canonical states verifies the conjugation
relations, $c^2=1$, and the absence of a fixed state on the complete
color-diagonal union.
\end{proof}
